% Auto-generated by scripts/exp_fold_gauge_anomaly_ldp_rate.py
\begin{proposition}[规范差的大偏差率函数接口（Legendre 变换 + 端点闭式）]\label{prop:fold-gauge-anomaly-ldp}
在均匀基线 $\omega\sim\mathrm{Unif}(\Omega_{m+1})$ 下，令 $G_m=\sum_{t=1}^{m}g_t$，其中 $g_t=\mathbf 1\{X_t\neq Y_t\}$ 为规范差的逐位失配指示。则对任意 $a\in[0,1]$，比例 $G_m/m$ 满足大偏差原理，其速率函数可由压力函数做 Legendre--Fenchel 变换给出：
\[
I(a)=\sup_{\theta\in\mathbb R}\bigl(a\theta-P_G(\theta)\bigr),\qquad P_G(\theta)=\log\rho(A_\theta),
\]
其中 $A_\theta$ 为命题 \ref{prop:fold-gauge-anomaly-pressure} 的 $4\times 4$ 加权邻接矩阵。并且端点的稀有事件指数具有完全闭式：
\[
\boxed{\ I(0)=\log(2/\varphi),\qquad I(1)=\log 2\ }.
\]
\end{proposition}

\begin{proof}[证明要点（可审计）]
由于 $g_t$ 是有限状态边移位上的可加观测量，G\"artner--Ellis 定理给出 $I$ 为 $P_G$ 的 Legendre--Fenchel 变换。端点闭式来自极限扭曲：当 $\theta\to-\infty$（即 $u=e^\theta\to 0$）时，失配边权趋于 $0$，谱半径退化为“全匹配子图”的 Perron 根，给出 $P_G(-\infty)=\log(\varphi/2)$，故 $I(0)=-P_G(-\infty)=\log(2/\varphi)$；当 $\theta\to+\infty$ 时，谱半径由全失配三环支配，满足 $P_G(\theta)=\theta-\log 2+o(1)$，故 $I(1)=\log 2$。
\end{proof}
